\chapter{Considerações Parciais}
\label{cap:consideracoes_parciais}


\section{Síntese parcial da pesquisa}
\label{sec:sintese_parcial}

Esta dissertação investiga, no nível da qualificação, como combinar sinais operacionais de máquina e evidências estatísticas históricas de rompimentos para inferir a origem de eventos de rompimento em uma fábrica de fios esmaltados de cobre e alumínio. A pergunta foi enunciada na Seção~\ref{sec:problema_pesquisa} e desdobrada nos objetivos específicos da Subseção~\ref{subsec:objetivos_especificos}.

Até este ponto, foram construídos os elementos fundacionais da proposta. A fundamentação do domínio industrial situa o problema do rompimento de fio nas tradições de qualidade e confiabilidade da manufatura contínua, e a revisão sistemática da literatura identifica a lacuna que a proposta enfrenta: a integração de detecção de anomalias, análise estatística histórica, fusão de evidências e explicabilidade em um único sistema aplicado à fabricação de fios esmaltados.

A arquitetura do WireSense foi formalizada no Capítulo~\ref{cap:proposta}, com dois pilares independentes, pipeline de dados e estrutura de saída de cada pilar. O Pilar de Sinais Operacionais opera por máquina, com \emph{Isolation Forest} treinado sobre intervalos sem incidente. O Pilar de Histórico Estatístico usa hora-máquina como métrica de exposição, Taxa Relativa por Material como indicador padronizado e cruzamento de falhas por máquina como sinal complementar.

Os elementos ainda em construção até a defesa completam a proposta com validação retrospectiva, implantação em regime de inferência continuada, painel de acompanhamento e discussão dos resultados. Esta qualificação não traz resultados experimentais, que constituem objeto do trabalho previsto entre esta etapa e a defesa.


\section{Cronograma para a defesa}
\label{sec:cronograma}

O cronograma de trabalhos previstos entre a qualificação e a defesa está organizado em seis etapas sequenciais, apresentadas na Tabela~\ref{tab:cronograma}. As etapas combinam atividades de implementação, experimentação, implantação e redação, e cobrem o período subsequente à qualificação até a defesa prevista para fevereiro de 2027.

\begin{table}[htbp]
    \centering
    \footnotesize
    \caption{Cronograma de trabalhos previstos até a defesa.}
    \label{tab:cronograma}
    \begin{tabularx}{\textwidth}{Xl}
        \toprule
        \textbf{Etapa}                                                                                 & \textbf{Prazo previsto} \\
        \midrule
        Refinamento da implementação dos dois pilares e da regra de integração                          & Jul--Ago 2026 \\
        Validação retrospectiva sobre o lote-caso confirmatório e demais eventos documentados           & Set--Out 2026 \\
        Implantação do gatilho de inferência continuada integrado ao sistema de chão de fábrica         & Nov 2026 \\
        Desenvolvimento do painel de acompanhamento e da camada de alertas                              & Dez 2026 \\
        Redação do capítulo de Resultados e Discussão                                                   & Jan 2027 \\
        Revisão geral do texto e preparação da defesa                                                   & Fev 2027 \\
        \bottomrule
    \end{tabularx}
    \fonte{Elaborado pelo autor (\the\year).}
\end{table}

A primeira etapa avança a implementação dos pilares, incorporando ajustes de engenharia e finalização da regra de integração. A segunda conduz a validação retrospectiva sobre dados históricos da fábrica, partindo de um lote de cobre com causa documentada por inspeção pericial e depois se estendendo a eventos documentados de devoluções de cliente e paradas internas. A expectativa é cobrir as classes de resposta previstas pela regra de integração apresentada na Seção~\ref{sec:regras_integracao}.

A terceira e a quarta etapas tratam da implantação em regime de inferência continuada, com gatilho acoplado ao sistema da máquina, retorno da classificação à interface operacional, painel de acompanhamento e camada de alertas. As duas últimas etapas tratam da redação final, incluindo a expansão deste capítulo em Resultados e Discussão e conclusão consolidada.


\section{Trabalhos futuros}
\label{sec:trabalhos_futuros}

Para além da defesa desta dissertação, identificam-se direções de extensão da pesquisa que ampliam o alcance do método proposto e o aproximam da operação cotidiana da fábrica.

A direção principal de extensão é a passagem do regime reativo para um regime preditivo. A implantação descrita nesta dissertação opera por gatilho, acionando o método de inferência quando o operador registra a parada por rompimento no sistema da máquina. A extensão preditiva inverte esse fluxo: a partir do monitoramento contínuo dos sensores e do histórico estatístico do lote em curso, o sistema antecipa a probabilidade de rompimento iminente e emite alerta preventivo. Essa direção exige tratamento adicional de calibração, definição de métricas de antecipação e gestão de falsos positivos, e constitui agenda de trabalho posterior à conclusão do mestrado.

A ampliação da cobertura de sensores e de tipos de defeito é uma segunda direção de extensão. O Pilar de Sinais Operacionais foi desenhado sobre um subconjunto de variáveis priorizadas, e o Pilar de Histórico Estatístico opera sobre os rompimentos. A incorporação de defeitos de bolha, asperidade e instabilidade de forno como categorias adicionais de evento e a inclusão de variáveis sensoriais hoje fora do conjunto selecionado permitiriam estender a inferência da origem para outros sintomas operacionais relevantes.

A inclusão do fornecedor como dimensão estatística direta dos indicadores do Pilar de Histórico Estatístico, complementar à preservação atual como atributo contextual descrita na Seção~\ref{sec:impl_historico}, é uma terceira direção de extensão. A construção de indicadores agregados por fornecedor permitiria identificar padrões sistêmicos de qualidade que se manifestam ao longo de múltiplos lotes da mesma origem, oferecendo evidência adicional para apoiar decisões de qualificação de fornecedores em horizonte estratégico, além do diagnóstico evento a evento que o WireSense produz.

A incorporação de realimentação de especialistas como fonte de rótulo fraco é uma quarta direção. Quando a equipe técnica de qualidade revisa um evento e atribui causa provável, esse julgamento pode ser registrado e usado como rótulo de baixa frequência mas alta confiança, refinando ao longo do tempo a calibração dos pilares e da regra de integração.

A generalização da arquitetura para outros processos de manufatura contínua com falhas multifatoriais é a quinta direção. A separação em pilares independentes, com fusão por regra determinística, é em princípio agnóstica ao domínio. Aplicações em fundição, laminação, extrusão e em outros processos com características semelhantes ao da fábrica analisada poderiam validar a transferibilidade do método proposto, com adaptação dos critérios de seleção de variáveis sensoriais e da estrutura categórica de matérias-primas.
